\section{Background and Related Work}
\label{sec:background_related}

\subsection{Delegation and the pre-authorization surface}

EIP-7702 introduces a transaction type containing signed authorization tuples. Processing a tuple writes a delegation indicator, $\mathtt{0xef0100}\,\|\,\textit{address}$, into the authority account; code-loading and execution follow the referenced address while preserving the authority account's execution context~\cite{eip7702}. The indicator is consequently a pointer-like designator, not the delegate program. Our input is the referenced runtime bytecode. This distinction also separates the point of intervention from post-hoc monitoring: the scorer is intended to inform a signer before authorization, while assuming only that an integration can supply or retrieve the runtime. It does not parse the authorization or perform the retrieval itself.

The closest EIP-7702 security study, Huang et al., characterizes EOA-targeted, contract-targeted, and composite attacks and combines historical transaction analysis with cross-contract static analysis~\cite{huang2026darkside}. Its released detection artifact uses Gigahorse-derived intermediate facts and Datalog rules. That study establishes the source of our positive labels, but its objective and information boundary differ from ours: it analyzes observed chain activity and reconstructed program facts, whereas we evaluate bytecode-only scoring before future transaction history is available. We did not execute the full USENIX Gigahorse/Datalog pipeline. Our sensitive-name rule approximation and external-call structural over-approximation are intentionally local bytecode baselines; neither reproduces that pipeline, and their performance cannot establish superiority to it.

Broader Ethereum phishing work covers adjacent surfaces. PTXPhish studies malicious transaction payloads and detects patterns using transaction and contract context~\cite{chen2025ptxphish}. PhishingHook instead applies multiple learned representations directly to EVM bytecode and opcodes for phishing-contract classification~\cite{derosa2025phishinghook}. ContractWard learns opcode-bigram models for smart-contract vulnerability classes~\cite{wang2021contractward}. These bytecode methods show that static learned representations can support screening without transaction histories, but their labels concern generic phishing or vulnerabilities rather than EIP-7702 delegate authorization. AuthGuard likewise uses bytecode only, but its target, input hygiene, family split, and pre-authorization decision boundary are specific to delegation.

These information boundaries determine which comparisons are meaningful. A pre-authorization bytecode score can operate without prior activity but observes only the candidate program; a post-hoc transaction detector can incorporate counterparties, ordering, value flow, and realized effects. Neither input subsumes the other. Accordingly, we compare learned and rule-based bytecode methods only inside the same task-aligned protocol, and use transaction-dependent studies to position the deployment stage rather than as metric baselines.

\subsection{Transaction-dependent phishing models}

A separate literature treats addresses as nodes and transfers or interactions as edges. Chen et al. use handcrafted account and temporal transaction features for Ethereum phishing detection~\cite{chen2020ijcai}, while Chen et al. model accounts and transactions with a graph convolutional network and autoencoder~\cite{chen2020toit}. Such address- and graph-level systems can exploit behavior distributed across counterparties and time. They are therefore complementary to bytecode scoring, not interchangeable baselines: their information may be absent for a fresh delegate, and their typical output is an address-level or transaction-level judgment after observable activity. Conversely, bytecode-only scoring cannot use signer-specific portfolios, social context, cash-out paths, or future behavior and should not be presented as replacing transaction analysis.

\subsection{Dependence, drift, and adversarial programs}

Security classifiers are especially sensitive to experimental dependence. TESSERACT identifies spatial and temporal biases in malware evaluation and notes that related variants across training and testing can inflate apparent performance~\cite{pendlebury2019tesseract}. Transcend addresses a different issue---detecting concept drift as deployed malware models age~\cite{jordaney2017transcend}. We borrow the former's evaluation discipline by keeping related bytecode in frozen families; we do not claim that family grouping eliminates all dependence, yields attacker attribution, or solves temporal drift. The random diagnostic instead quantifies how much the split choice changes results on this corpus.

Adversarial-malware research further shows why feature-space perturbations are insufficient for executable artifacts. Grosse et al. adapt adversarial examples to constrained malware features~\cite{grosse2017adversarial}, and Pierazzi et al. formalize problem-space attacks in which transformations must produce realizable programs while accounting for side-effect features~\cite{pierazzi2020problemspace}. Our mutations follow this problem-space motivation but make a narrower guarantee: metadata, embedded-address, selector, and dead-code changes are structure-preserving transformations under our opcode-skeleton checker. We neither prove formal EVM equivalence nor evaluate arbitrary recompilation or dynamic control-flow obfuscation.

Finally, static analyzers recover richer program properties than a feature scorer. Gigahorse decompiles EVM bytecode into an intermediate representation suitable for declarative analyses~\cite{grech2019gigahorse}; Securify derives dependency information and checks compliance and violation patterns~\cite{tsankov2018securify}. These systems illustrate heavyweight semantic and control-flow analysis, while AuthGuard performs lightweight deterministic feature extraction and prediction. The approaches answer different questions and may be composed. We did not benchmark either analyzer, so we make no speedup claim and do not infer that heavyweight analysis is infeasible at the decision point.
